\section{Method}
\label{sec:method}

\subsection{Overview and Notation}
\label{sec:overview}

Let $f_\theta$ be a frozen pure-text LLM with $L$ layers and hidden width
$d$, and let $\mathcal{S}\subset\{1,\dots,L\}$ be the set of layers we
modulate. An image is processed into $N$ patches by the teacher's image
processor. \modelname{} never places those patches inside the context of the
deeper layers. Instead, it produces a compact set of cell features
$V\in\mathbb{R}^{N_c\times d_v}$ and uses them to rotate the inputs of the
seven linear projections ($q,k,v,o,\text{gate},\text{up},\text{down}$) in
layers $\mathcal{S}$, as illustrated in Fig.~\ref{fig:arch}.

Two model variants share this skeleton. \emph{\modelname-Tower} takes $V$
from a frozen external vision encoder and is the direct successor of
rotation-based injection. \emph{\modelname-Free} additionally deletes the
encoder: $V=\hat{v}$ is produced by the student's own first $K$ layers, and
all image tokens are dropped afterwards. Both variants are followed by the
same on-policy distillation stage.

\subsection{Tier 1: Built-in Vision Frontend}
\label{sec:frontend}

\paragraph{Static rotation with channel scaling.}
A pure orthogonal rotation preserves norms and angles but cannot change
singular values, so it can re-aim features without re-weighting channels. For a
linear projection $y = x W^{\top} + b$, we first rotate each $b$-dimensional
input block and then apply a learned per-channel diagonal scaling
$d\in\mathbb{R}^{d_{\text{in}}}$ before multiplying by $W^{\top}$, i.e.\
$y = ((x R_{\mathrm{s}})\odot d)W^{\top} + b$, which is equivalent to the
transformed weight
\begin{equation}
W' = W\,\mathrm{diag}(d)\,R_{\mathrm{s}}^{\top},
\qquad
R_{\mathrm{s}} = (I-A)^{-1}(I+A),
\label{eq:static}
\end{equation}
where $A$ is skew-symmetric, $R_{\mathrm{s}}\in SO(b)$ acts on $b{=}64$
channel blocks, and $d\in\mathbb{R}^{d_{\text{in}}}$ is initialised to ones.
Because $R_{\mathrm{s}}$ is orthogonal, training stays well conditioned and
the rescaling supplies the degrees of freedom that orthogonality removes. For
$K{=}4$ layers and seven projections on a $1.7$B backbone, the static
in-layer module has $2.40$M parameters (versus $2.49$M for rank-$16$ LoRA and
$201$M for duplicating the four layers), while all encoder-free variants share
the same $3.67$M patch/positional embedding and $50.36$M $2{\times}2$ cell
merger ($56.42$M vs.\ $56.52$M total frontend parameters).

\paragraph{Reading the image once.}
Image patches are linearly embedded and given row and column positional
embeddings. They are then propagated through the student's first $K$ layers
with \emph{bidirectional} attention under Eq.~\eqref{eq:static}, which lets
every patch see the whole image---the regime a vision encoder needs and the
regime a causal LLM never provides. Adjacent $2\times2$ patches are
concatenated and projected by a two-layer merger into one cell feature,
\begin{equation}
\hat{v} = \mathrm{MLP}\big(\mathrm{LN}([h_{4i},h_{4i+1},h_{4i+2},h_{4i+3}])\big)
\in\mathbb{R}^{d_v},
\label{eq:merge}
\end{equation}
yielding $\hat{v}\in\mathbb{R}^{N_c\times d_v}$ with $N_c=N/4$.

\paragraph{Discard, then modulate.}
After Eq.~\eqref{eq:merge} the $N$ image patch tokens are deleted: the text
sequence of length $T$ is then processed through all $L$ layers of the frozen
LLM without any visual tokens in its self-attention sequence (paying
$O(K N^2 + L T^2)$ attention instead of $O(L(T+N_c)^2)$). The only trace of the
image during the text pass is $\hat{v}\in\mathbb{R}^{N_c\times d_v}$, which
drives the deep modulation of Sec.~\ref{sec:dynamic}. This separates our
two-phase design from built-in-encoder approaches that keep visual tokens in the
sequence across all $L$ layers.

\paragraph{Optional feature supervision.}
When a teacher encoder is available during training we align $\hat{v}$ with
its cell features $v^{\ast}$ through
\begin{equation}
\mathcal{L}_{\mathrm{feat}} = \frac{1}{N_c}\sum_{i}\Big[(1-\cos(\hat v_i,v^\ast_i))
+ \lambda_1\,\ell_{\mathrm{smooth}}\big(\tfrac{\hat v_i}{s},\tfrac{v^\ast_i}{s}\big)\Big],
\label{eq:feat}
\end{equation}
with $s=\mathrm{std}(v^\ast)$ and $\lambda_1{=}0.5$. This term only shapes the
frontend during training; it is not needed at inference and, as we show in
Sec.~\ref{sec:abl}, the model still learns to read images when it is
switched off entirely.

\subsection{Tier 2: Dynamic Rotation Injection}
\label{sec:dynamic}

\paragraph{Generating a rotation per layer.}
For a layer $\ell\in\mathcal{S}$ we keep $N_b$ learnable queries
$Q_\ell\in\mathbb{R}^{N_b\times d_q}$ ($d_q{=}256$). Optionally, under
\emph{question conditioning} (\texttt{qcond}), the mean-pooled prompt hidden
state $\bar x_{\text{prompt}}$ at layer $\ell$ is projected by a
zero-initialised linear layer $W_{qc}^{(\ell)}\mathrm{LN}(\bar x_{\text{prompt}})$
and added to $Q_\ell$ so the rotation queries depend on the question as well as
the image. Visual cell features $V\in\mathbb{R}^{N_c\times d_v}$ are
projected to keys and values, and the queries read out a context
\begin{equation}
C_\ell = \mathrm{MHA}\big(Q_\ell,\; W_k^{(\ell)}V,\; W_v^{(\ell)}V\big),
\qquad
\theta_\ell = W_o^{(\ell)}C_\ell ,
\label{eq:ctx}
\end{equation}
where $\theta_\ell\in\mathbb{R}^{N_b\times b(b-1)/2}$ collects the strictly
lower-triangular entries of the skew-symmetric blocks. The rotation is
\begin{equation}
A_r = \mathrm{tril}(\theta_{\ell,r}) - \mathrm{tril}(\theta_{\ell,r})^{\top},
\qquad
R_{\ell,r} = (I-A_r)^{-1}(I+A_r),
\label{eq:cayley}
\end{equation}
and a rotated projection computes
\begin{equation}
y = \big(x\,R_{\ell}\big)W^{\top}+b,
\qquad
R_\ell = \mathrm{blkdiag}(R_{\ell,1},\dots,R_{\ell,N_b}),
\label{eq:rotlin}
\end{equation}
with $x$ split into $N_b$ channel blocks. All seven projections of a layer
share the same visual context, which is computed once and cached for the whole
forward pass, so generation pays for the rotation generator once per layer
rather than once per step.

\paragraph{Multi-frequency rotation (\texttt{mf}).}
A single pooled context forces every block to look at the image through the
same aperture. We split the $N_b$ blocks into three bands, following the
frequency division of rotary position embeddings~\citep{su2024roformer}. 
Low-frequency blocks attend to the globally pooled feature $\bar v$;
mid-frequency blocks
attend to all cells with a smooth positional encoding
($\mathrm{base}{=}10^4$); high-frequency blocks use a rapidly varying encoding
($\mathrm{base}{=}50$, softmax temperature $\tau{=}0.25$) to sharpen
localisation. A zero-initialised gate then mixes rotation with identity,
\begin{equation}
\tilde x_r = x_r + 2\sigma\big(W_g x\big)_r\,\big(x_rR_{\ell,r}-x_r\big),
\label{eq:gate}
\end{equation}
so the model starts from the pretrained function and adds rotation only where
it helps. This variant favours multi-step reasoning over diagrams, where the
model must repeatedly consult a small region.

\paragraph{Spatial-anchor and direct narrow-subspace hybrid (\texttt{hyb} / \texttt{direct\_hyb}).}
The second variant keeps the orthogonal rotation and adds a narrow-subspace
cross-attention residual before the frozen linear projection,
\begin{equation}
\Delta x = W_{\mathrm{out}}\,
\mathrm{MHA}\big(W_q\,\mathrm{LN}(x),\, K_a,\, V_a\big),
\qquad
y=\big(xR_\ell+\Delta x\big)W^{\top}+b,
\label{eq:hyb}
\end{equation}
where $W_{\mathrm{out}}$ is zero-initialised. In the pooled-anchor variant
(\texttt{hyb}), $K_a,V_a\in\mathbb{R}^{M\times d_q}$ come from $M{=}16$
learnable spatial anchors pooled from $V$; in the uncompressed direct variant
(\texttt{direct\_hyb}), $K_a,V_a\in\mathbb{R}^{N_c\times d_q}$ project all $N_c$
visual cells directly into the $d_q{=}256$ narrow subspace (with a post-attention
LayerNorm before $W_{\mathrm{out}}$). In both cases the cross-attention operates
in $d_q{=}256\ll d$ and no visual token enters the LLM's self-attention KV
cache. To isolate the contribution of orthogonal rotation from this additive
cross-attention path, we also evaluate a strict \emph{cross-attention-only}
ablation (\texttt{no\_rot}, $R_\ell\equiv I$, removing the rotation generator
entirely).

All variants are controlled by explicit flags while the LLM backbone stays
strictly frozen.

\subsection{Modality Decoupling}
\label{sec:decouple}

The rotations act through a context-local wrapper that is active only inside
the image branch. When no visual context is set, each wrapped projection calls
its original weight, so a text-only input follows exactly the computation of
the unmodified LLM and produces bit-identical logits. This is what allows
\modelname{} to claim lossless text ability while still being trained on image
data. It also explains a negative result we report in
Sec.~\ref{sec:eff}: any variant that merges vision adapters into the text
pathway---LoRA, static rotation, or rotation with scaling---loses most of its
arithmetic ability, which is why the decoupled design is not a convenience but
a requirement.

\subsection{Step 2: Online Policy Distillation}
\label{sec:opd}

Teacher-forced training only visits prefixes of ground-truth answers. Greedy
decoding visits the student's own prefixes, where a small early error moves
the hidden state off the training manifold; rotation-injected models are
particularly sensitive because the rotation is conditioned on the current
hidden state itself. We therefore distil on-policy.

Let $\hat y_{1:T}\sim\pi_S(\cdot\mid I,q)$ be a rollout sampled from the
student with the vision context active. A frozen teacher VLM
$\pi_T$ (an instruction-tuned model of the same tokenizer family) scores the
\emph{same} trajectory in one forward pass, giving
\begin{equation}
\mathcal{L}_{\mathrm{OPD}}=
\mathbb{E}_{\hat y}\Big[\sum_{t=1}^{T}
\mathrm{KL}\!\left(\pi_T(\cdot\mid I,q,\hat y_{<t})\,\Big\|\,
\pi_S(\cdot\mid I,q,\hat y_{<t})\right)
+\lambda_{\mathrm{ce}}\,\mathcal{L}_{\mathrm{CE}}\Big],
\label{eq:opd}
\end{equation}
where the cross-entropy term anchors the model to gold tokens that appear in
the rollout. For \modelname-Free{} we add Eq.~\eqref{eq:feat} with a small
weight to keep $\hat v$ stable while the frontend is updated. Training uses
$500$ steps, batch size $4$, learning rates $2\times10^{-5}$ (rotation) and
$10^{-5}$ (frontend), and takes under fifteen minutes on one GPU. The
distillation is thus a cheap post-hoc stage rather than a second training run.

\subsection{Complexity}
\label{sec:complexity}

For a text length $T$ and $N$ visual patches, a token-concatenation layer
attends over $T+N$ positions, costing $O((T+N)^2d)$; the visual part alone is
$O(2TNd+N^2d)$ and is paid in every layer. \modelname{} pays $O(N^2d)$ for
$K$ early layers in the frontend and, in the remaining $L-K$ layers, only
$O(T^2d)$ for self-attention plus $O(TN d_q)$ for the anchor residual, with
$d_q\ll d$. The block rotations add $O(Tdb)$, independent of both $N$ and the
number of projections' output width in the leading term. Measured at $4096$
patches the total drops from $9.85$ to $2.54$ TFLOPs
(Tab.~\ref{tab:eff}), and because the frontend encodes patches once with
bidirectional attention, wall-clock latency at high resolution improves as
well.
